% answers_short.tex —— 2016 年·浙江卷·理科 · 简短答案速查
% 由 tools/make_short_answers.py 自 answers.tex 的 \gksolhead/\gkanswer 提取，
% 并与 2016/2016_数学真题_浙江理_答案要点.md 逐题比对后生成，勿手改（改源 MD 后重跑）。

\begin{gknotice}
\textbf{使用说明}：本节只给最终答案便于快速核对，完整推导见解析版。答案为编写者独立推导，\textbf{非官方标准答案}。
\end{gknotice}

\gkpart{选择题}
\begin{center}\begin{tabular}{ccccccccc}
\toprule
题号 & 1 & 2 & 3 & 4 & 5 & 6 & 7 & 8 \\
\midrule
答案 & B & C & C & D & B & A & A & D \\
\bottomrule
\end{tabular}\end{center}

\gkpart{填空题}
\begin{itemize}[leftmargin=2.6em]
  \item[\textbf{9.}] \(9\)
  \item[\textbf{10.}] \(A=\sqrt{2}\)，\(b=1\)
  \item[\textbf{11.}] 表面积 \(72\,\mathrm{cm}^{2}\)，体积 \(32\,\mathrm{cm}^{3}\)
  \item[\textbf{12.}] \(a=4\)，\(b=2\)
  \item[\textbf{13.}] \(a_{1}=1\)，\(S_{5}=121\)
  \item[\textbf{14.}] \(\dfrac{1}{2}\)
  \item[\textbf{15.}] \(\dfrac{1}{2}\)
\end{itemize}

\gkpart{解答题}
\begin{itemize}[leftmargin=2.6em]
  \item[\textbf{16.}] (1) 见解析版；(2) \(A=\dfrac{\pi}{4}\) 或 \(\dfrac{\pi}{2}\)（两解均满足题设）
  \item[\textbf{17.}] (1) 见解析版；(2) 二面角 \(B-AD-F\) 的平面角的余弦值为 \(\dfrac{\sqrt{3}}{4}\)
  \item[\textbf{18.}] (1) \(x\in[2,2a]\)（当 \(x\geqslant 1\) 时化为 \((x-2)(x-2a)\leqslant 0\)；当 \(x<1\) 时无解）；(2) ① 当 \(3\leqslant a\leqslant 2+\sqrt{2}\) 时 \(m(a)=0\)，当 \(a>2+\sqrt{2}\) 时 \(m(a)=4a-a^{2}-2\)（因 \(F(1)=0\) 恒成立，最小值取 \(\min\{0,\,4a-a^{2}-2\}\)）；② 当 \(3\leqslant a\leqslant 4\) 时 \(M(a)=34-8a\)，当 \(a>4\) 时 \(M(a)=2\)
  \item[\textbf{19.}] (1) 截得的线段长为 \(\dfrac{2a^{2}|k|\sqrt{1+k^{2}}}{1+a^{2}k^{2}}\)；(2) 离心率的取值范围是 \(\left(0,\dfrac{\sqrt{2}}{2}\right]\)
  \item[\textbf{20.}] (1) 见解析版；(2) 见解析版
\end{itemize}
